\section{引言}

「明朝是一个复杂系统，它的灭亡是涌现出来的」——这句话今天很常见，也很少被算过。
涌现（emergence）在复杂网络研究里不是一句形容：它要求一个系统的宏观形状
既不等于任何单个成员的性质，又是从成员之间的连接规则里长出来的。
把这句话搬到一部叙事抽取出的知识图谱上，它就变成一组可算的量：
度数分布有多集中、分布是不是幂律、图是不是一整块、连接的形状随时间怎么移。
本篇要做的不是复述"兴衰=涌现"，而是先把它拆成可复核的四个量，
再报告其中哪些算得出来、哪些在现有样本量下必须收手。

具体问法是四个问题。\textbf{第一，这部叙事写的人有多集中：}
1231 位人物里，有多少人一条关系也没有，度数的分布形状是平台还是长尾。
\textbf{第二，这个形状能不能叫无标度：}
按 Clauset 等人对幂律拟合的要求，尾指数应当在大样本上是稳的；
若它随截断点的抬升系统性漂移，那么"幂律"这一标签就不该贴上去。
\textbf{第三，图是不是连通的：}
一个 dynastic 叙事的网络通常被渲染成一张密网，但密网是可以由少数枢纽撑起来的假象。
\textbf{第四，这个形状有没有随时间换位：}
若"崩塌"真的在书写中被感知为一个过程，事件类别的比重应当在三期之间迁移，
而不只是总量变化。

语料是通俗明史叙事《明朝那些事儿》全七部，以及从同一语料抽取、可独立重建的
结构化知识图谱：人物 1231、事件 1106、关系 2324 条、地点 563 处。
其中 1006 件事件带可解析年份，落在 1328—1645 之间的 218 个年份上。
本篇量的不是明朝是否真亡于某个原因，而是这部被读过最多的叙事，
在结构上长成了什么样的一张图。

结果指向一个清楚的形状。度数侧高度集中但不必称幂律：
基尼系数 0.8199，537 人（43.6\%）一条关系也没有，
度数排序前一百的 108 人（占全部人物的 8.8\%）握有全部度数和的 65.1\%，
尾指数估计值却随 xmin 从 1.806 一路漂到 2.777（漂移 0.971，尾样本仅 142 条）。
连通侧是一整块：741 个顶点带边、18 个连通块、最大块 703 点，
占带边顶点的 94.9\%，只占全部人物的 57.1\%。
时间侧不是均匀增长：年均 4.615 件、变异系数 1.16、按年基尼 0.5069，
而 1368—1499 / 1500—1599 / 1600—1662 三期之间的事件类别分布
χ²=129.05（df=20，p 小于 $10^{-6}$），Cramér's V=0.259。

贡献有三点。
\begin{itemize}
\item 把"兴衰是涌现"拆成四个可复算的量：集中（基尼）、分布形状（尾指数稳定性）、
连通（最大连通块占比）、时间迁移（类别占比的三期位移），每一个都给出点数的口径。
\item 明确一条负结果——本篇不主张幂律，并把"在多大的尾样本上必须放弃幂律主张"
写成一条可复核的判据（第 6 节），而不是留作脚注。
\item 报一个可用卡方复核的时期迁移：军事战争占比 28.9 / 21.8 / 40.9，
政治斗争占比 33.9 / 40.8 / 21.5，两者在第三期方向相反。
\end{itemize}

边界先说死：本篇不主张这套网络形状对应明代社会的真实结构，
也不主张"崩塌"二字落在哪一年。它量的是一部叙事在抽取之后长成的形状。

第 2 节回顾相关文献，第 3 节交代语料与口径，第 4 节写方法，
第 5 节报告结果，第 6 节讨论效度边界与幂律的样本量上限，第 7 节收束。
